%% =====================================================================
%%  CAPITULO 3 - TRABALHOS RELACIONADOS
%% =====================================================================

\chapter{Trabalhos Relacionados}
\label{cap:relacionados}

Este capítulo posiciona o trabalho em relação a três linhas: processos
orientados por especificação, agentes de linguagem aplicados ao desenvolvimento
de software e verificação de software assistida por agentes. A comparação é
predominantemente qualitativa; onde uma comparação quantitativa exigiria dados
que não foram consultados na fonte original, isso é declarado.

\section{Processos orientados por especificação}
\label{sec:rel-sdd}

O desenvolvimento orientado por testes estabelece o teste como artefato de
primeira classe, escrito antes da implementação (\citeonline{beck2002tdd}). A
especificação por exemplo leva a ideia adiante, substituindo descrição abstrata
por exemplos concretos que funcionam simultaneamente como requisito, teste e
documentação (\citeonline{adzic2011}). Ambos compartilham o mecanismo central
adotado neste trabalho: o critério de aceitação precede a implementação.

A engenharia de requisitos normatizada fornece o vocabulário para tornar o
critério verificável, incluindo atributos de qualidade de requisito
(\citeonline{iso29148}). O modelo de qualidade de produto
(\citeonline{iso25010}) fornece a taxonomia para classificar requisitos não
funcionais.

O desenvolvimento de linguagens específicas de domínio ilustra o caso extremo da
mesma tese: quando o domínio é bem compreendido, uma notação restrita reduz a
ambiguidade e a possibilidade de erro (\citeonline{mernike2005}). O paralelo com
este trabalho está na função dos identificadores estáveis: ao nomear cada
requisito, critério e tarefa, restringe-se o espaço de referências ambíguas.

\section{Agentes de linguagem no desenvolvimento de software}
\label{sec:rel-agentes}

Três trabalhos representativos organizam múltiplos agentes para produzir software.
MetaGPT atribui papéis análogos aos de uma equipe e representa artefatos
intermediários de forma estruturada, com o argumento de que representações
intermediárias bem definidas reduzem a propagação de erro entre etapas
(\citeonline{hong2024metagpt}). ChatDev modela o desenvolvimento como
conversação entre papéis de projetista, programador e revisor, com foco na
comunicação entre agentes como mecanismo de correção
(\citeonline{qian2024chatdev}). AutoGen trata a conversação entre agentes
configuráveis como primitiva de construção, permitindo compor fluxos com
participação humana (\citeonline{wu2023autogen}).

Esses trabalhos demonstram que a decomposição em papéis melhora a qualidade do
artefato gerado. Entretanto, o critério de sucesso predominante é a qualidade do
código ou a taxa de conclusão da tarefa, e não a existência de uma cadeia
auditável entre requisito, teste e evidência. Em particular, nenhum dos três
define condições de parada por informação ausente como parte do protocolo: o
agente pode continuar produzindo na ausência de dados essenciais.

\section{Padrões de raciocínio e autocorreção}
\label{sec:rel-padroes}

O intercalamento entre raciocínio e ação permite que o agente use o resultado de
uma ação para decidir a próxima, corrigindo a trajetória durante a execução
(\citeonline{yao2023react}). A reflexão sobre o próprio erro, com registro da
avaliação crítica reutilizada em tentativas posteriores, melhora o desempenho em
tarefas que admitem tentativa e erro (\citeonline{shinn2023reflexion}). A
explicitação do raciocínio intermediário melhora o desempenho em tarefas de
múltiplos passos (\citeonline{wei2022}).

Esses padrões são adotados no framework proposto, mas com uma restrição: a
autocorreção é aplicada a artefatos internos (especificação, tarefas, código),
nunca à evidência. A evidência de teste não é objeto de reflexão do agente; ela é
resultado de execução. Essa separação é deliberada e responde ao problema de
opacidade apontado na Seção~\ref{sec:agentes}.

\section{Comparação}
\label{sec:rel-comparacao}

O Quadro~\ref{qua:comparacao} compara os trabalhos segundo cinco critérios
derivados do problema de pesquisa.

\begin{quadro}[htb]
	\caption{Comparação entre abordagens correlatas segundo critérios de controle de processo}
	\label{qua:comparacao}
	\centering
	\footnotesize
	\begin{tabularx}{\textwidth}{@{}p{3.3cm} c c c c c@{}}
		\toprule
		\textbf{Abordagem} &
		\rotatebox{70}{\textbf{Especificação prévia}} &
		\rotatebox{70}{\textbf{Múltiplos papéis}} &
		\rotatebox{70}{\textbf{Critério verificável}} &
		\rotatebox{70}{\textbf{Evidência de execução}} &
		\rotatebox{70}{\textbf{Parada por dado ausente}} \\
		\midrule
		\citeonline{beck2002tdd}      & Parcial & Não  & Sim     & Sim     & Não \\
		\addlinespace
		\citeonline{adzic2011}        & Sim     & Não  & Sim     & Sim     & Não \\
		\addlinespace
		\citeonline{iso29148}         & Sim     & Não  & Sim     & Não     & Não \\
		\addlinespace
		\citeonline{hong2024metagpt}  & Parcial & Sim  & Não     & Não     & Não \\
		\addlinespace
		\citeonline{qian2024chatdev}  & Parcial & Sim  & Não     & Não     & Não \\
		\addlinespace
		\citeonline{wu2023autogen}    & Não     & Sim  & Não     & Parcial & Não \\
		\addlinespace
		\citeonline{yao2023react}     & Não     & Não  & Não     & Sim     & Não \\
		\addlinespace
		\citeonline{shinn2023reflexion} & Não   & Não  & Não     & Sim     & Não \\
		\addlinespace
		\midrule
		\textbf{Processo proposto}    & Sim     & Sim  & Sim     & Sim     & Sim \\
		\bottomrule
	\end{tabularx}
	\legend{Fonte: elaborado pelo autor (2026). A classificação é qualitativa e baseada na leitura
	dos trabalhos; ``parcial'' indica que o critério existe de forma implícita ou incompleta.
	Confirmar cada classificação contra o texto original --- ver P-16 a P-21.}
\end{quadro}

\section{Posicionamento e contribuição pretendida}
\label{sec:rel-posicionamento}

A comparação do Quadro~\ref{qua:comparacao} evidencia que as abordagens existentes
cobrem bem a geração cooperativa e os padrões de raciocínio, mas não combinam, em
um único processo, quatro elementos: especificação prévia obrigatória, critério
verificável com identificador estável, evidência de execução como única fonte de
conclusão e parada explícita diante de informação ausente.

A contribuição pretendida não é um novo modelo de linguagem nem um novo algoritmo
de coordenação. É um \emph{protocolo de controle} aplicável sobre agentes
existentes, com três propriedades verificáveis:

1. \textbf{Rastreabilidade} --- cada entrega é referenciável a um requisito.
2. \textbf{Falseabilidade} --- cada requisito relevante tem um teste capaz de reprová-lo.
3. \textbf{Honestidade epistêmica} --- ausência de dado é registrada como ausência, nunca suprida por estimativa.

A avaliação dessa contribuição é feita por estudo de caso, conforme o
Capítulo~\ref{cap:metodos}, e não por experimento controlado. Essa escolha
metodológica tem custo e benefício, discutidos na Seção~\ref{sec:ameacas}.
